\documentclass{beamer}
\usetheme{breadhouse}

\title{A minimal example for bread house}
\date{\today}
\author{Bread House Owner}
\institute{House for Bread}
\begin{document}
\maketitle

\section{Introduction}

\begin{frame}{Introduction}

\begin{itemize}
  \item Your introduction goes here!
  \item Use \texttt{itemize} to organize your main points.
\end{itemize}

\vskip 1cm

\begin{block}{Examples}
Some examples of commonly used commands and features are included, to help you get started.
\end{block}

\end{frame}

\section{Some \LaTeX{} Examples}

\subsection{Tables and Figures}

\begin{frame}{Tables and Figures}

\begin{itemize}
\item Use \texttt{tabular} for basic tables --- see Table~\ref{tab:widgets}, for example.
\item You can upload a figure (JPEG, PNG or PDF) using the files menu.
\item To include it in your document, use the \texttt{includegraphics} command (see the comment below in the source code).
\end{itemize}

 Commands to include a figure:
\begin{figure}
 \includegraphics[width=.1\textwidth]{logo.eps}
 \caption{\label{fig:your-figure}Caption goes here.}
\end{figure}

\begin{table}
\centering
\begin{tabular}{l|r}
Item & Quantity \\\hline
Widgets & 42 \\
Gadgets & 13
\end{tabular}
\caption{\label{tab:widgets}An example table.}
\end{table}

\end{frame}

\section{Mathematics}

\begin{frame}{Readable Mathematics}
Let $X_1, X_2, \ldots, X_n$ be a sequence of independent and identically distributed random variables.
\end{frame}
\end{document}
